\section{Results}
\subsection{Replication Results}

Table~\ref{tab:side_replication_full} compares the original ordered logistic regression results with the reproduced SIDE results. For each human-judgment dimension, we keep the original paper metrics and append a new row, \textit{SIDE (new)}, which reports the reproduced SIDE result from our local replication.

\begin{table*}[t]
\centering
\caption{Original ordered logistic regression results and reproduced SIDE results. The row \textit{SIDE (new)} reports the reproduced result from this replication.}
\label{tab:side_replication_full}
\small
\begin{tabular}{llrrrrr}
\hline
\textbf{Dimension} & \textbf{Metric} & \textbf{OR} & \textbf{Value} & \textbf{Std. Error} & \textbf{$t$-value} & \textbf{$p$-value} \\
\hline

\multirow{11}{*}{Overall DA Score}
& BLEU-1 & 0.9990 & -0.0010 & 0.0017 & -0.5822 & 0.6222 \\
& BERTScore-R & 1.0029 & 0.0029 & 0.0023 & 1.2634 & 0.2943 \\
& SentenceBERT-CS & 1.0057 & 0.0057 & 0.0020 & 2.8882 & 0.0100 \\
& InferSent-CS & 0.9980 & -0.0020 & 0.0025 & -0.8064 & 0.5250 \\
& ROUGE-1-P & 1.0058 & 0.0058 & 0.0018 & 3.2552 & 0.0033 \\
& ROUGE-4-R & 1.0024 & 0.0024 & 0.0011 & 2.1987 & 0.04667 \\
& ROUGE-W-R & 0.9996 & -0.0004 & 0.0017 & -0.2485 & 0.8040 \\
& c\_coeff & 1.0143 & 0.0142 & 0.0012 & 11.4549 & $<0.0001$ \\
& CodeT5-plus-CS & 1.0044 & 0.0044 & 0.0017 & 2.5383 & 0.0220 \\
& SIDE & 1.0205 & 0.0203 & 0.0017 & 11.7029 & $<0.0001$ \\
& \textbf{SIDE (new)} & \textbf{1.0194} & \textbf{0.0192} & \textbf{0.0018} & \textbf{10.8722} & \textbf{$<0.0001$} \\
\hline

\multirow{11}{*}{Content Adequacy}
& BLEU-1 & 0.9672 & -0.0333 & 0.0347 & -0.9615 & 0.4200 \\
& BERTScore-R & 1.0525 & 0.0512 & 0.0468 & 1.0927 & 0.3929 \\
& SentenceBERT-CS & 1.1121 & 0.1063 & 0.0406 & 2.6172 & 0.0225 \\
& InferSent-CS & 1.0000 & 0.0000 & 0.0517 & -0.0004 & 1.0000 \\
& ROUGE-1-P & 1.0943 & 0.0901 & 0.0364 & 2.4733 & 0.0260 \\
& ROUGE-4-R & 0.9914 & -0.0086 & 0.0226 & -0.3814 & 0.7811 \\
& ROUGE-W-R & 1.0437 & 0.0428 & 0.0351 & 1.2202 & 0.3700 \\
& c\_coeff & 1.3729 & 0.3169 & 0.0255 & 12.4350 & $<0.0001$ \\
& CodeT5-plus-CS & 1.1413 & 0.1322 & 0.0358 & 3.6877 & $<0.0001$ \\
& SIDE & 1.6265 & 0.4864 & 0.0366 & 13.2942 & $<0.0001$ \\
& \textbf{SIDE (new)} & \textbf{1.5739} & \textbf{0.4535} & \textbf{0.0373} & \textbf{12.1635} & \textbf{$<0.0001$} \\
\hline

\multirow{11}{*}{Conciseness}
& BLEU-1 & 1.0656 & 0.0636 & 0.0342 & 1.8576 & 0.1575 \\
& BERTScore-R & 1.0160 & 0.0159 & 0.0470 & 0.3375 & 0.8770 \\
& SentenceBERT-CS & 1.0587 & 0.0571 & 0.0403 & 1.4163 & 0.2617 \\
& InferSent-CS & 1.0079 & 0.0079 & 0.0511 & 0.1547 & 0.8770 \\
& ROUGE-1-P & 1.0075 & 0.0075 & 0.0359 & 0.2091 & 0.8770 \\
& ROUGE-4-R & 0.9937 & -0.0063 & 0.0228 & -0.2766 & 0.8770 \\
& ROUGE-W-R & 1.0954 & 0.0911 & 0.0346 & 2.6352 & 0.0267 \\
& c\_coeff & 1.2433 & 0.2178 & 0.0254 & 8.5798 & $<0.0001$ \\
& CodeT5-plus-CS & 1.0630 & 0.0611 & 0.0355 & 1.7212 & 0.1700 \\
& SIDE & 1.3844 & 0.3253 & 0.0354 & 9.1835 & $<0.0001$ \\
& \textbf{SIDE (new)} & \textbf{1.3462} & \textbf{0.2973} & \textbf{0.0361} & \textbf{8.2382} & \textbf{$<0.0001$} \\
\hline

\multirow{11}{*}{Fluency}
& BLEU-1 & 1.0650 & 0.0629 & 0.0345 & 1.8243 & 0.2267 \\
& BERTScore-R & 1.0559 & 0.0544 & 0.0467 & 1.2038 & 0.4600 \\
& SentenceBERT-CS & 0.9958 & -0.0042 & 0.0404 & -0.1043 & 0.9170 \\
& InferSent-CS & 1.0543 & 0.0529 & 0.0515 & 1.0272 & 0.4600 \\
& ROUGE-1-P & 1.0334 & 0.0329 & 0.0365 & 0.9004 & 0.4600 \\
& ROUGE-4-R & 1.0304 & 0.0299 & 0.0229 & 1.3061 & 0.4600 \\
& ROUGE-W-R & 0.9672 & -0.0334 & 0.0347 & -0.9603 & 0.4600 \\
& c\_coeff & 1.2080 & 0.1553 & 0.0253 & 6.1437 & $<0.0001$ \\
& CodeT5-plus-CS & 1.0249 & 0.0246 & 0.0357 & 0.6902 & 0.5444 \\
& SIDE & 1.2826 & 0.2489 & 0.0359 & 6.9359 & $<0.0001$ \\
& \textbf{SIDE (new)} & \textbf{1.2647} & \textbf{0.2348} & \textbf{0.0366} & \textbf{6.4227} & \textbf{$<0.0001$} \\
\hline
\end{tabular}
\end{table*}

The reproduced SIDE results are highly consistent with the original paper. Across all four human-judgment dimensions, \textit{SIDE (new)} remains statistically significant with \(p<0.0001\), and every reproduced odds ratio is greater than 1. This means that higher reproduced SIDE scores are still associated with higher odds of receiving better human ratings.

For Overall DA Score, the reproduced SIDE odds ratio is 1.0194, very close to the original SIDE value of 1.0205. This is one of the strongest pieces of evidence for a successful replication because the overall score is the broadest human quality dimension. The reproduced coefficient is slightly smaller than the original coefficient, but the direction, magnitude, and significance remain almost unchanged.

For Content Adequacy, \textit{SIDE (new)} obtains an odds ratio of 1.5739. This is lower than the original SIDE value of 1.6265, but it is still much larger than the other metrics in the original table. The closest non-SIDE metric is c\_coeff with an odds ratio of 1.3729, followed by CodeT5-plus-CS with 1.1413 and SentenceBERT-CS with 1.1121. This confirms the central role of SIDE: it is especially strong at predicting whether a generated summary contains the right semantic content for the code.

For Conciseness, the reproduced SIDE odds ratio is 1.3462, again slightly below the original SIDE value of 1.3844 but still higher than all other metrics in the original table. The next strongest metric is c\_coeff with an odds ratio of 1.2433, while most lexical and embedding metrics are close to 1 or statistically insignificant. This suggests that SIDE is not only measuring content overlap; summaries that are better aligned with code semantics are also more likely to be judged concise.

For Fluency, \textit{SIDE (new)} has an odds ratio of 1.2647, compared with the original SIDE value of 1.2826. This dimension is less directly tied to code semantics than content adequacy, yet SIDE still outperforms the other metrics in the table. The closest competing metric is c\_coeff with an odds ratio of 1.2080, while CodeT5-plus-CS and most reference-overlap metrics are weaker. This indicates that SIDE captures a quality signal that is broader than lexical similarity alone.

The reproduction preserves the ranking pattern of the original study. SIDE remains the strongest metric among the reported candidates for Content Adequacy, Conciseness, and Fluency, and it remains highly competitive for Overall DA Score. The small decreases from the original SIDE row to \textit{SIDE (new)} are expected in a local artifact rerun because of environment differences, dependency versions, and recovered data assets. However, these differences do not change the interpretation: the reproduced SIDE metric continues to show a robust and statistically significant association with human judgments, especially for semantic content quality.


\subsection{Extension Results}

The Python extension evaluates whether SIDE filtering improves downstream code summarization training. Table~\ref{tab:extension_metrics} compares the unfiltered control condition with the SIDE-filtered condition; Table~\ref{tab:output_diagnostics} reports output-length diagnostics and empty-summary rates. The results are organized around two questions:

\begin{itemize}

    \item Does SIDE filtering improve average evaluation scores and strengthen the correlation between SIDE and other quality signals?

    \item Does SIDE filtering reduce unstable generation behavior, such as empty summaries or abnormal output lengths?
\end{itemize}

\begin{table*}[t]
\centering
\caption{Comparison between unfiltered and SIDE-filtered training conditions. Correlation values are Spearman correlations with SIDE.}
\label{tab:extension_metrics}
\begin{tabular}{lrrrrrr}
\hline
\textbf{Metric} &
\textbf{No Filter Avg.} &
\textbf{SIDE-Filtered Avg.} &
\textbf{$\Delta$ Avg.} &
\textbf{No Filter $\rho$} &
\textbf{SIDE-Filtered $\rho$} &
\textbf{$\Delta \rho$} \\
\hline
BLEU-4 & 5.4121 & 6.3185 & +0.9064 & 0.50 & 0.54 & +0.04 \\
ROUGE-L & 33.4254 & 34.7324 & +1.3070 & 0.55 & 0.60 & +0.05 \\
METEOR & 17.0302 & 18.5237 & +1.4935 & 0.53 & 0.57 & +0.04 \\
ChrF & 24.3055 & 27.4066 & +3.1011 & 0.55 & 0.57 & +0.02 \\
TF-IDF Cosine & 17.1617 & 17.5747 & +0.4130 & 0.52 & 0.59 & +0.07 \\
GPT-5.3 Judge & 2.5080 & 3.1020 & +0.5940 & 0.48 & 0.54 & +0.06 \\
\hline
\end{tabular}
\end{table*}

\paragraph{\textbf{Can filtered training data improve summarization quality?}} Table~\ref{tab:extension_metrics} shows that SIDE filtering improves all average evaluation scores and also strengthens the monotonic relationship between SIDE and the other metrics. This indicates that SIDE filtering has two effects: it improves the generated summaries according to standard evaluation metrics, and it makes SIDE more consistent with the other quality signals.

For lexical overlap, BLEU-4 increases from 5.4121 to 6.3185, a gain of 0.9064 points. This suggests that the SIDE-filtered model produces summaries with more exact phrase-level overlap with the references. Although BLEU-4 is strict and sensitive to wording, the improvement shows that semantic filtering does not hurt surface-level matching; instead, it slightly improves it. The Spearman correlation between BLEU-4 and SIDE also rises from 0.50 to 0.54, meaning that summaries with higher SIDE scores are more consistently ranked higher by BLEU-4 after filtering.

ROUGE-L increases from 33.4254 to 34.7324, a gain of 1.3070 points. Since ROUGE-L measures longest common subsequence overlap, this improvement suggests that SIDE-filtered training helps the model preserve more reference-relevant content in its generated summaries. Its correlation with SIDE also improves from 0.55 to 0.60, which is the strongest post-filtering correlation among the reference-based metrics. This shows that the filtered model produces summaries whose sequence-level overlap is more aligned with semantic code-summary alignment.

METEOR increases from 17.0302 to 18.5237, a gain of 1.4935 points. This is important because METEOR is more tolerant of flexible wording than BLEU and rewards unigram alignment, stemming, and paraphrase-like similarity. The improvement suggests that SIDE filtering helps the model generate summaries that are semantically closer to the reference, even when exact wording differs. The correlation with SIDE also increases from 0.53 to 0.57.

ChrF shows the largest automatic-metric gain, increasing from 24.3055 to 27.4066, a gain of 3.1011 points. Because ChrF is based on character n-gram overlap, this improvement indicates that the SIDE-filtered model produces outputs with stronger fine-grained lexical similarity to the reference. Its correlation with SIDE increases from 0.55 to 0.57, showing a smaller but still positive improvement in metric alignment.

TF-IDF cosine increases from 17.1617 to 17.5747, a smaller gain of 0.4130 points. Although the average improvement is modest, the correlation with SIDE rises from 0.52 to 0.59, which is the largest correlation gain in the table. This suggests that while TF-IDF similarity changes only slightly in average value, its ranking behavior becomes much more consistent with SIDE after filtering.

The GPT-5.3 judge score increases from 2.5080 to 3.1020, a gain of 0.5940 points on a 1--5 scale. This is the additional evidence that the improvement is not limited to automatic reference-based metrics. The LLM judge evaluates whether a summary is faithful and useful for the code, so this increase suggests that SIDE-filtered training improves perceived summary quality. Its correlation with SIDE also rises from 0.48 to 0.54, indicating that the judge becomes more aligned with SIDE under the filtered condition.

We can see every average score increases and every Spearman correlation with SIDE also increases from Table~\ref{tab:extension_metrics}. The average-score gains show that SIDE filtering improves summary quality, while the correlation gains show that the filtered model produces outputs whose quality rankings are more coherent across metrics. This supports the main extension claim that SIDE filtering is not merely removing data; it is selecting training examples that lead to more stable and semantically aligned summarization behavior.


\begin{table}[t]
\centering
\caption{Output length diagnostics for the 500 generated summaries in each condition. Token counts are computed after removing sequence boundary markers.}
\label{tab:output_diagnostics}
\small
\begin{tabular}{lrrr}
\hline
\textbf{Diagnostic} & \textbf{No Filter} & \textbf{SIDE-Filtered} & \textbf{$\Delta$} \\
\hline
Mean length & 7.934 & 9.398 & +1.464 \\
25th percentile & 6 & 6 & 0 \\
Median length & 7 & 8 & +1 \\
75th percentile & 9 & 11 & +2 \\
Minimum length & 0 & 0 & 0 \\
Maximum length & 42 & 34 & -8 \\
Empty summaries & 23 & 6 & -17 \\
Empty summary rate & 4.60\% & 1.20\% & -3.40 pp \\
\hline
\end{tabular}
\end{table}

\paragraph{\textbf{How do the output diagnostics differ after training with different datasets?}} Table~\ref{tab:output_diagnostics} shows that SIDE filtering also improves output stability. The mean generated-summary length increases from 7.934 tokens to 9.398 tokens. More importantly, the distribution shifts upward: the median increases from 7 to 8 tokens, and the 75th percentile increases from 9 to 11 tokens. This suggests that the SIDE-filtered model is not only producing a few longer summaries; its typical output is also more informative.

The lower quartile remains 6 tokens in both conditions, so SIDE filtering does not eliminate short summaries entirely. However, it substantially reduces the most severe failure case: empty outputs. The unfiltered model produces 23 empty summaries out of 500, corresponding to an empty-summary rate of 4.60\%. The SIDE-filtered model produces only 6 empty summaries, reducing the empty-summary rate to 1.20\%. This is a practical improvement because empty summaries are unusable regardless of their metric scores.

The maximum length also decreases from 42 to 34 tokens. Combined with the higher median and higher upper quartile, this suggests that SIDE filtering makes the output distribution more stable: it reduces empty generations while avoiding extremely long outputs. Thus, the output diagnostics support the metric results. SIDE filtering improves several quality measures, increases agreement between SIDE and other metrics, and reduces degenerate generation behavior.
